% P16-C5: finite counterfactual semantics and the potential-outcome bridge.
\section{Finite counterfactual reasoning and potential outcomes}
\label{sec:p16-finite-counterfactuals}

Pearl defines a unit's counterfactual by solving an equation-replaced
model with that unit's same background attributes
\cite[Definition~4 and Eq.~(27), pp.~120--121]{Pearl2009CausalOverview}.
The earlier presentation uses the same structural interpretation and
consistency relation \cite[author's rejoinder, p.~703]{Pearl1995CausalDiagrams}.
Here the finite construction keeps one common outcome space for every
world. These initial constructions permit arbitrarily dependent
disturbances; the acyclic graph and total structural functions supply
the required unique solutions.

\subsection{Several worlds for the same unit}

\begin{hbdefinition}{def:p16-counterfactual-worlds}{Counterfactual variables and their common background}
Fix the finite model \(\mathcal M\), its outcome space
\((\Omega,q)\), and its solution \(V\) from
\cref{def:p16-finite-scm,thm:p16-causal-solution}.
For \(I\subseteq N\), \(x\in S_I\), and \(J\subseteq N\),
the counterfactual tuple
\[
 V_J[I=x](\omega)\defeq V_J^{I=x}(\omega)
\]
is the solution on coordinates \(J\) after the equations at \(I\)
are replaced, evaluated at the same \(\omega\). The bracket notation
keeps the intervention label explicit. For a named variable \(Y\),
write \(Y[I=x]\) for its coordinate; Pearl's \(Y_x(u)\) has
this interpretation when \(u\) is the background disturbance vector.

For an integer \(m\geq0\), let \(\mathcal W=\{0,\ldots,m\}\) be a finite set of world
labels, with interventions \((I_k,x^k)\) and queried coordinate
sets \(J_k\). Define the random tuple
\[
 Q(\omega)=\bigl(V_{J_k}[I_k=x^k](\omega)\bigr)_{k\in\mathcal W}.
\]
Its \emph{joint cross-world law} is the ordinary distribution of
this tuple on the one probability space. Incompatible interventions
at different labels are evaluated in different modified models;
they do not impose inconsistent equations within one model.
Taking \(I_0=\varnothing\) includes the factual world.

For two successive interventions, write
\(V[I=x;J=y]\) for evaluation after first replacing \(I\)'s
equations and then replacing \(J\)'s equations. If the sets overlap,
the second replacement is the operative equation on their intersection.
\index{counterfactual!common background}\index{cross-world distribution}
\end{hbdefinition}

\begin{hbtheorem}{thm:p16-counterfactual-laws}{proved=annex:p16.counterfactual-laws}{Effectiveness, consistency, and valid composition}
For \(I,J\subseteq N\), \(x\in S_I\), \(y\in S_J\), and
every original outcome \(\omega\):
\begin{enumerate}
\item \(V_I[I=x](\omega)=x\) (effectiveness).
\item If \(V_I(\omega)=x\), then
      \(V[I=x](\omega)=V(\omega)\) (factual consistency).
\item If \(V_J[I=x](\omega)=y\), then
      \(V[I=x;J=y](\omega)=V[I=x](\omega)\)
      (composition at attained values).
\item Successive replacement equals replacement on \(I\cup J\)
      with the assignment equal to \(y\) on \(J\) and to \(x\)
      on \(I\setminus J\). The replacements commute when their
      assignments agree on \(I\cap J\), in particular when the
      sets are disjoint.
\item If \(j\notin I\cup\operatorname{de}_G(I)\), then
      \(V_j[I=x](\omega)=V_j(\omega)\).
\end{enumerate}
Composition uses values attained in the already modified world.
Substituting values from a different intervention world does not follow
from this assertion. All five conclusions hold also at zero-mass outcomes.
\end{hbtheorem}

\begin{hbdefinition}{def:p16-counterfactual-evidence}{Positive evidence and the three-stage calculation}
Let \(E\subseteq\Omega\) be specified evidence with
\(\mathbb P(E)>0\), for example \(E=\{V_B=b\}\) evaluated
in the factual world. Define
\[
 q_E(\omega)=\frac{q(\omega)\mathbf1_E(\omega)}{\mathbb P(E)}.
\]
The \emph{abduction/action/prediction calculation} consists of:
\begin{enumerate}
\item \emph{Abduction}: compute the posterior background weights
      \(q_E\) using the specified evidence event.
\item \emph{Action}: replace the equations for each queried world,
      retaining these weights and the same underlying outcome.
\item \emph{Prediction}: evaluate each world in a topological order
      and sum the posterior weights for the desired joint result.
\end{enumerate}
The evidence event is held fixed during this calculation. Replacing
it by the analogous event in an intervened world specifies another
query. If \(\mathbb P(E)=0\), this conditional calculation is
undefined under the finite conditioning convention.
\index{counterfactual!abduction action prediction}
\end{hbdefinition}

\begin{hbtheorem}{thm:p16-counterfactual-calculation}{proved=annex:p16.counterfactual-calculation}{The finite joint law and conditional calculation}
The cross-world tuple \(Q\) has a normalized mass function. For
any set \(D\) of its possible tuple values and positive evidence \(E\),
\[
 \mathbb P(Q\in D\mid E)
     =\sum_{\omega:Q(\omega)\in D}q_E(\omega).
\]
The three-stage procedure terminates and returns this probability.
Each unconditional world marginal is its previously defined intervention law.

Let \(U=(U_i)_{i\in N}\), with finite value set
\(\mathcal U=\prod_i\mathcal U_i\). Every queried world depends
on \(\omega\) only through \(U(\omega)\). Consequently the
calculation may group the posterior weights into
\[
 r_E(u)=\sum_{\omega:U(\omega)=u}q_E(\omega).
\]
The joint posterior \(r_E\) must be retained; independence of the
original disturbances does not imply their independence after evidence.
\end{hbtheorem}

\subsection{A finite twin construction}

\begin{hbdefinition}{def:p16-twin-data}{Shared background and duplicated endogenous coordinates}
For the finite world family above, create a single new root \(R\)
with value set \(\mathcal U\) and equation \(R=U\) on the
original outcome space. For every \(k\in\mathcal W\) and
\(i\in N\), create a distinct vertex \((i,k)\), with variable
\(T_i^k\) taking values in \(S_i\). Define
\[
 T_i^k=
 \begin{cases}
 x_i^k,&i\in I_k,\\
 f_i\bigl((T_j^k)_{j\in\operatorname{pa}_G(i)},R_i\bigr),
       &i\notin I_k,
 \end{cases}
\]
where \(R_i\) is the \(i\)-th component of \(R\).
An unmodified copied equation has its within-world parent edges
and an edge from \(R\). A constant copied equation has no parents.
Use the disturbance \(U\) for the shared root and constant
disturbances for all copied equations. For two worlds this is a
\emph{twin construction}; the same definition handles any finite family.
\end{hbdefinition}

\begin{hbtheorem}{thm:p16-finite-twin-construction}{proved=annex:p16.finite-twin-construction}{The twin model represents the exact joint counterfactual law}
The constructed graph is finite and acyclic, and
\[
 T_i^k(\omega)=V_i[I_k=x^k](\omega)
\]
for every \(i,k,\omega\). Its marginal joint law on the queried
copies, and its conditional law given any retained positive evidence
event, therefore coincide with the corresponding law of \(Q\).

If \(r(u)=\mathbb P(U=u)\), the mass of a complete copied
assignment \((v^k)_k\) is
\[
 \sum_{u\in\mathcal U}r(u)
 \prod_{k\in\mathcal W}
 \left[
 \prod_{i\in I_k}\mathbf1_{\{v_i^k=x_i^k\}}
 \prod_{i\notin I_k}
 \mathbf1_{\{v_i^k=f_i(v^k_{\operatorname{pa}_G(i)},u_i)\}}
 \right].
\]
The root represents the full joint background law even when its
components are dependent. Drawing an independent background separately
for each copied world would generally give a different joint law.
\end{hbtheorem}

\subsection{Binary-treatment potential outcomes}

The source's consistency and conditional-exchangeability formulas appear
in Eqs.~(36)--(40) \cite[pp.~127--129]{Pearl2009CausalOverview}.
The results below give explicit finite hypotheses and their proofs.

\begin{hbdefinition}{def:p16-binary-potential-outcomes}{Potential outcomes, consistency, and target populations}
For a binary treatment vertex \(t\) and a distinct response vertex
\(r\) with finite real value set, write
\[
 A=V_t\in\{0,1\},\quad Y=V_r,\quad
 Y(0)=V_r[t=0],\quad Y(1)=V_r[t=1].
\]
The source's \(Y_a\) is written \(Y(a)\) here. Let \(L=V_C\)
be a finite covariate tuple with
\(C\cap(\{t,r\}\cup\operatorname{de}_G(\{t\}))=\varnothing\).
It is unaffected by these treatment interventions.

An abstract finite potential-outcome model instead specifies jointly
\((A,L,Y(0),Y(1))\) on a finite probability space and defines
the observed response by \emph{consistency},
\[
 Y=Y(A)=(1-A)Y(0)+AY(1).
\]
Define the individual contrast \(\Delta=Y(1)-Y(0)\), the
\emph{average treatment effect}
\(\operatorname{ATE}=\mathbb E[\Delta]\), and, when
\(\mathbb P(A=1)>0\), the \emph{average treatment effect among
the treated}
\(\operatorname{ATT}=\mathbb E[\Delta\mid A=1]\).
For binary responses the probability of benefit is
\(\mathbb P(Y(0)=0,Y(1)=1)\).

The condition \(Y(a)\mathrel{\perp\!\!\!\perp}A\mid L\)
means ordinary conditional independence on every positive covariate
stratum, using the common joint law that includes the potential outcome.
It is a causal assumption about that law, not a conclusion from observing
\((A,L,Y)\). No independence between \(Y(0)\) and \(Y(1)\)
is imposed.
\end{hbdefinition}

\begin{hbtheorem}{thm:p16-binary-po-bridge}{proved=annex:p16.binary-po-bridge}{The exact finite binary-treatment bridge}
The structural definitions satisfy consistency pointwise and
\[
 \mathbb P(Y(a)=y)
  =\mathbb P_{\mathcal M}(Y=y\mid\operatorname{do}(A=a)).
\]
Conversely, every finite joint potential-outcome model satisfying the
displayed consistency relation is represented by a finite acyclic
structural model with a shared latent root, the same joint
\((A,L,Y(0),Y(1),Y)\) law, and these same treatment-intervention
responses. This converse concerns the single binary treatment and
its baseline covariates; it asserts no general counterfactual
axiomatization for arbitrary cyclic systems.
\end{hbtheorem}

\begin{hbcorollary}{thm:p16-back-door-exchangeability}{proved=annex:p16.back-door-exchangeability}{A structural sufficient condition for exchangeability}
If the original structural model has mutually independent disturbances
and \(C\) satisfies the back-door criterion for the treatment and
response, then \(Y(a)\mathrel{\perp\!\!\!\perp}A\mid L\)
for each \(a=0,1\). This conclusion itself needs no treatment
overlap; using it to identify conditional ratios requires the
positivity assumptions below.
\end{hbcorollary}

\begin{hbtheorem}{thm:p16-po-ate-identification}{proved=annex:p16.po-ate-identification}{Identification of marginal potential outcomes and ATE}
In the finite potential-outcome model assume consistency, and, for
each \(a=0,1\), assume
\(Y(a)\mathrel{\perp\!\!\!\perp}A\mid L\).
Require \(\mathbb P(A=a,L=\ell)>0\) for both \(a\) whenever
\(\mathbb P(L=\ell)>0\). Put
\(\mu_a(\ell)=\mathbb E[Y\mid A=a,L=\ell]\) on these strata.
Then
\[
\begin{aligned}
 \mathbb P(Y(a)=y)
  &=\sum_{\ell:\mathbb P(L=\ell)>0}
       \mathbb P(Y=y\mid A=a,L=\ell)\mathbb P(L=\ell),\\
 \operatorname{ATE}
  &=\sum_{\ell:\mathbb P(L=\ell)>0}
       [\mu_1(\ell)-\mu_0(\ell)]\mathbb P(L=\ell).
\end{aligned}
\]
These quantities are identified from the observed law of \((A,L,Y)\)
within the class satisfying these assumptions. Joint exchangeability
\((Y(0),Y(1))\mathrel{\perp\!\!\!\perp}A\mid L\) is not
assumed. The joint coupling need not be identified, as
\cref{thm:p16-po-coupling-nonidentification} demonstrates.
\end{hbtheorem}

\begin{hbtheorem}{thm:p16-po-att-identification}{proved=annex:p16.po-att-identification}{Identification of ATT under control-outcome exchangeability}
Assume consistency, \(\mathbb P(A=1)>0\), and only
\(Y(0)\mathrel{\perp\!\!\!\perp}A\mid L\).
Require \(\mathbb P(A=0,L=\ell)>0\) for every \(\ell\)
with \(\mathbb P(A=1,L=\ell)>0\). Define \(\mu_0,\mu_1\)
on these treated strata by the same observed conditional means. Then
\[
 \operatorname{ATT}
 =\sum_{\ell:\mathbb P(A=1,L=\ell)>0}
     [\mu_1(\ell)-\mu_0(\ell)]\mathbb P(L=\ell\mid A=1).
\]
This identifies ATT from the observed law. Exchangeability of \(Y(1)\)
is not required for this conclusion, and no treated observations are
required in covariate strata absent from the treated population.
\end{hbtheorem}

\subsection{A complete posterior and treatment-effect calculation}

\begin{hbdefinition}{def:p16-cf-worked-data}{Covariate-dependent treatment assignment}
Take independent \(U_L\), a fair bit, and \(U_A\), uniform on
\(\{1,2,3,4\}\), with equations
\[
 L=U_L,\qquad A=\mathbf1_{\{U_A\leq1+2L\}},\qquad
 Y=L+A(1+L).
\]
The response disturbance is constant. The graph is
\(L\to A\), \(L\to Y\), \(A\to Y\), with response values
in \(\{0,1,3\}\). The eight outcomes \((u_L,u_A)\) have
equal mass \(1/8\). All three endogenous variables are observed.
\end{hbdefinition}

\begin{hbtheorem}{thm:p16-cf-worked-calculation}{proved=annex:p16.cf-worked-calculation}{Abduction, action, prediction, ATE and ATT in one model}
The model's complete observed law is
\[
\begin{array}{ccc|c}
 L&A&Y&\text{mass}\\\hline
 0&0&0&3/8\\
 0&1&1&1/8\\
 1&0&1&1/8\\
 1&1&3&3/8
\end{array}
\]
It satisfies the ATE and ATT assumptions and has
\[
 \operatorname{ATE}=3/2,\qquad\operatorname{ATT}=7/4,
 \qquad\mathbb E[Y\mid A=1]-\mathbb E[Y\mid A=0]=9/4.
\]
For factual evidence \(E=\{A=1,Y=3\}\), the posterior gives
mass \(1/3\) to each of \((1,1),(1,2),(1,3)\).
After replacing \(A\)'s equation by zero, each of those outcomes
has response one, so
\[
 \mathbb P(Y(0)=1\mid E)=1.
\]
The population intervention probability \(\mathbb P(Y(0)=1)\)
is \(1/2\), and the observed conditional \(\mathbb P(Y=1\mid A=0)\)
is \(1/4\). These are different queries with different weightings.
\end{hbtheorem}

\subsection{What randomized treatment does not identify}

\begin{hbtheorem}{thm:p16-po-coupling-nonidentification}{proved=annex:p16.po-coupling-nonidentification}{Joint potential outcomes and probability of benefit}
There are two finite acyclic models with the same allowed graph
\(A\to Y\), randomized fair treatment independent of the entire
potential-outcome pair, and identical full observational and
single-world intervention laws on \((A,Y)\), but different couplings:
\[
\begin{array}{cc|cc}
 Y(0)&Y(1)&\text{model }\mathcal M_{=}&\text{model }\mathcal M_{\ne}\\\hline
 0&0&1/2&0\\
 0&1&0&1/2\\
 1&0&0&1/2\\
 1&1&1/2&0
\end{array}
\]
Both have \(\mathbb P(Y(a)=1)=1/2\) at either treatment and
ATE and ATT zero. Their probabilities of benefit are respectively
zero and \(1/2\). Even their responder-specific probabilities
\(\mathbb P(Y(0)=0\mid A=1,Y=1)\) differ, being zero and one.

More generally, for binary potential outcomes with marginal success
probabilities \(p_0,p_1\), the probability of benefit \(b\)
has the sharp bounds
\[
 \max(0,p_1-p_0)\leq b\leq\min(p_1,1-p_0).
\]
Every value in this interval is realized by such a randomized finite
structural model with the given marginal potential-outcome laws.
\end{hbtheorem}

These results separate computation in a fully specified structural
model from identification in a class of models compatible with given
data. The intervention-identification algorithms have their stated
completeness domains in the following sections, including
\cref{thm:p16-tcc-completeness}.
Natural direct and indirect effects, their identifying assumptions
and their countermodels are supplied in
\cref{thm:p16-mediation-countable-formula,thm:p16-mediation-kernel-formula,
thm:p16-mediation-crossworld-countermodels,thm:p16-mediation-diagonal}.
General cyclic-model semantics and unrestricted counterfactual
axiomatizations retain their distinct requirements.
